\section{Définition de l'atlas généalogique et conventions d'évaluation métrologique}

Cet atlas n'est pas un tableau de coïncidences. Son noyau mathématique conserve
cinq champs inséparables ; le statut logique et le critère de réfutation sont énoncés ensuite, sans prendre la place d'aucun des cinq :

\begin{equation}
\boxed{
\begin{gathered}
\text{antécédent}\ \big|\ \text{équation}\ \big|\
\ \text{invariant conservé}\\
\text{signification}\ \big|\ \text{valeur terminale}
\end{gathered}}
\qquad
\begin{gathered}
[\text{statut};\\ \text{critère de réfutation}]
\end{gathered}.
\label{eq:atlas-six-fields}
\end{equation}


\section{Corollaire HMT de confluence métrologique}
\label{sec:metrological-confluence}

\begin{definition}[Quintuplet métrologique]
Soit \(x\) une grandeur obtenue par une réalisation HMT--MD. Son
\emph{quintuplet
métrologique} est
\[
 \mathfrak C(x)=
 \bigl(A_x\mid E_x\mid I_x\mid S_x\mid V_x\bigr),
\]
où \(A_x\) est l'antécédent typé, \(E_x\) l'équation de
détermination, \(I_x\) l'invariant conservé par la réalisation,
\(S_x\) la signification
structurelle et \(V_x\) la valeur terminale. Il n'est pas permis de former \(V_x\) avant
d'avoir fixé les quatre autres composantes. Deux sorties ayant la même valeur décimale
ne sont pas égales si leurs antécédents ou leurs invariants diffèrent.
\end{definition}

\begin{corollary}[Confluence métrologique]
\label{cor:metrological-confluence}
Les sorties de vide, d'action, d'information, de température, de gravité,
de criticité, de spectre, de mélange et de cosmologie construites dans les
\cref{ch:vacuum,ch:metrology,ch:chaos,ch:spectral,ch:gauss,ch:gravity,ch:modular,ch:ew,ch:cosmo}
admettent les quintuplets du \cref{tab:metrological-confluence}. Dans chaque
ligne, la valeur de la cinquième colonne est une conséquence des quatre premières, et
non une donnée qui les sélectionne. Les lignes partagent des nœuds généalogiques
---mode \(30\), \(\log3\), déterminant \(D_A\), canaux \(90/120\) et changement
\(6\to8\)---, mais conservent leurs domaines et leurs opérateurs séparés.
\end{corollary}

\begingroup
\setlength{\tabcolsep}{2pt}
\renewcommand{\arraystretch}{1.22}
\scriptsize
\begin{longtable}{@{}
>{\raggedright\arraybackslash}p{25mm}
>{\raggedright\arraybackslash}p{37mm}
>{\raggedright\arraybackslash}p{24mm}
>{\raggedright\arraybackslash}p{29mm}
>{\raggedright\arraybackslash}p{32mm}@{}}
\caption{Corollaire de confluence métrologique HMT : les cinq colonnes forment
une seule unité probatoire.}\label{tab:metrological-confluence}\\
\toprule
Antécédent & Ecuaci\'on & Invariant conservé & Significado estructural &
Valeur terminale\\
\midrule
\endfirsthead
\multicolumn{5}{l}{\small\itshape Suite du corollaire de
confluence métrologique HMT}\\
\toprule
Antécédent & Ecuaci\'on & Invariant conservé & Significado estructural &
Valeur terminale\\
\midrule
\endhead
\midrule
\multicolumn{5}{r}{\itshape Suite à la page suivante}\\
\endfoot
\bottomrule
\endlastfoot

\multicolumn{5}{@{}l}{\bfseries I. Action}\\*
deux réalisations de la droite d'action
(\cref{eq:action-twoway,eq:action-rapidity}) &
\(R_{\rm act}=H_5/\eta_{\rm ret}\),
\(\xi_{\rm act}=\tfrac12\log R_{\rm act}\) &
Décennie et orientation de feuillet &
L'action absolue et son quotient relatif correspondent à des fonctionnelles
distinctes ; la
décennie n'est pas utilisée comme ajustement &
\(R_{\rm act}\), \(\xi_{\rm act}\)\\

\addlinespace
\multicolumn{5}{@{}l}{\bfseries II. Vide électromagnétique}\\*
\(T=q_+P_++q_-P_-\), \(30\to90/120\)
(\cref{eq:vacuum-operator,eq:three-four-completion}) &
\(V=(I-T^{90})(I-T^{120})^{-1}\) &
Positivité, spectre des feuillets et complétion \(3\to4\) &
Le vide est un opérateur bicouche ; \(90=3\cdot30\) et \(120=4\cdot30\)
déterminent le quatrième secteur sans introduire de constante SI &
\(r_+=0.9999996285\ldots\), \(r_-=0.9997241371\ldots\)\\

réponse spectrale
(\cref{eq:vacuum-four,eq:vacuum-traces}) &
\(\widehat\varepsilon=r_+^2\),
\(\widehat\mu=r_-^2\),
\(\widehat Z=r_-/r_+\),
\(\widehat c=(r_+r_-)^{-1}\) &
Produit et quotient des feuillets, stables sous \(V\mapsto V\otimes I_{9^m}\) &
Permittivité, perméabilité, impédance et propagation sont quatre fonctions
du même opérateur positif &
\(0.9999992571\ldots\), \(0.9994483505\ldots\),
\(0.9997245085\ldots\), \(1.0002763105\ldots\)\\

droites constitutives
(\cref{eq:vacuum-sections}) &
\(\mu\varepsilon=c^{-2}\), \(\mu/\varepsilon=Z^2\) &
Homogénéité dans \(L_S,L_C,L_Q\) &
La carte dimensionnelle est une section ultérieure, et non une unité accolée à un
nombre adimensionnel &
\(Z,\mu,\varepsilon,c\) comme sections exactes\\

\((Z_{\rm vac},h_a,\alpha_{\HMT})\)
(\cref{eq:charge,eq:quantum-electrical-identities}) &
\(e_a^2=2\alpha h_a/Z\), \(R_KG_0=2\), \(K_J\Phi_0=1\) &
Produit bidirectionnel et orientation opposée du flux/de la fréquence &
Les quanta électriques descendent du vide et de l'action ; la charge SI
n'agit pas comme sélecteur &
\(e_a,R_K,G_0,\Phi_{0,a},K_{J,a}\)\\

\addlinespace
\multicolumn{5}{@{}l}{\bfseries III. Information, température et rayonnement}\\*
trit uniforme et structure chronologique \(108\)
(\cref{eq:boltzmann-clock,eq:log3-triple}) &
\(k_B\Theta_{\rm clk}\log3=h_{\rm int}/(108t_0)=E_P\) &
Entropie d'une décision ternaire et période périodicité temporelle de (108) phases &
Un même quotient mesure l'information, la température et l'énergie sans identifier
leurs trois espaces d'états &
\(\log3=1.0986122886681098\ldots\)\\

\(G\hbar\) et périodicité temporelle de (108) phases
(\cref{eq:planck-energy,eq:planck-trit}) &
\(E_P=\pi\hbar/(54t_0)=k_B\Theta_{\rm clk}\log3\) &
Produit bidirectionnel \(G\hbar\) et contenu informationnel TRIT &
Le quantum géométrique de Planck et le quantum informationnel associé à un trit
confluent dans une même section d'énergie &
\(E_P=h/(108t_0)\)\\

spectre de Planck
(\cref{eq:wien-lambda,eq:wien-frequency}) &
\(\lambda_{\max}/\ell_0=108\log3/x_5\),
\(\nu_{\max}t_0=x_3/(108\log3)\) &
Même spectre, jacobiens distincts &
Les maxima par longueur et par fréquence sont des réalisations fonctionnelles d'un antécédent commun, et non deux
approximations de la même coordonnée &
\(23.896756779\ldots\), \(0.0237794888\ldots\)\\

continu bosonique
(\cref{eq:stefan-hmt,eq:photon-number,eq:photon-entropy}) &
\(n_\gamma\ell_P^3=2\zeta(3)/[\pi^2(\log3)^3]\) et lois \(T^4,T^3\) &
Occupation bosonique et contenu informationnel par trit &
Radiance, nombre, énergie et entropie forment une famille thermique ; ce ne
sont pas quatre calibrations indépendantes &
Coefficients exacts dans \(\pi,\zeta(3),\log3\)\\

\addlinespace
\multicolumn{5}{@{}l}{\bfseries IV. Gravité, aire et mémoire modulaire}\\*
clôture circulaire périodicité temporelle de (108) phases
(\cref{thm:gravity-circular-selector,eq:gravity-theta108}) &
\(G=\theta_{108}c^5t_0^2/\hbar\),
\(\theta_{108}=(54/\pi)^2\) &
Égalité du rayon de Compton et du rayon gravitationnel &
La constante gravitationnelle est une section sélectionnée par une clôture de
longueurs, et non un nombre greffé &
\(G_{\rm int}=(54/\pi)^2c^5t_0^2/\hbar\)\\

\(A_5/C_5\to P_5\)
(\cref{thm:gravity-p5-coupling,eq:gravity-G5-spectrum}) &
\(\mathcal G_5=GP_5\), \(P_5^2=P_5\), \(\Tr P_5=5\) &
Rang central et réduction TT &
Cinq composantes du support tensoriel précèdent les deux hélicités
physiques ; elles ne se confondent pas avec elles &
\(\Spec\mathcal G_5=\{G^{(5)},0^{(7)}\}\)\\

changement \(6\to8\), dérivation de Barbero--Immirzi et canaux \(90/120\)
(\cref{eq:modular-area-operator,eq:modular-68-two-readers}) &
\(A_{\rm phys}/\ell_P^2=\gammaH a_0(N_{90}+N_{120})\) &
Quantum d'aire, orientation hexade--octade et \(\ell_P^2\) &
\(\gammaH\) fixe le quantum géométrique ; cette monographie adopte cette
grandeur sans rouvrir sa dérivation &
\(\gammaH a_0=0.0016755940749\ldots\)\\

occupation \(N\) et mémoire \(D\)
(\cref{thm:modular-area-memory,eq:modular-inverse-channels}) &
\(D=90N_{90}+120N_{120}\),
\(N_{90}=4N-D/30\), \(N_{120}=D/30-3N\) &
Déterminant \(30\) et provenance du canal &
L'aire quantifie l'occupation totale ; le hamiltonien holographique
modulaire conserve le secteur \(90/120\) de provenance &
\((N_{90},N_{120})\) reconstruit exactement\\

état de bord
(\cref{thm:modular-tfd,eq:modular-oriented-information,thm:modular-first-law}) &
\(I_{\rm or}=36\log3-S(\rho_A)\),
\(\delta S=\beta_{90}\delta\langle\mathcal A_{90}\rangle+
\beta_{120}\delta\langle\mathcal A_{120}\rangle\) &
Spectre de Schmidt et rapport \(\beta_{120}/\beta_{90}=4/3\) &
Entropie d'intrication, information orientée et aire sont trois
lectures compatibles du même bord &
\(I_{\rm or}=0.0016691832\ldots\),
\((\beta_{90},\beta_{120})=(5.97265\ldots,7.96353\ldots)\)\\

réseau local infini
(\cref{thm:modular-typeIII,eq:modular-typeIII-lambda}) &
\(\lambda_\partial=\exp(-30\Dmod)\) &
Groupe des rapports de poids sous persistance \(90/120\) &
La limite est un générateur modulaire GNS, et non une matrice de densité globale
de classe trace dans l'UHF &
\(\lambda_\partial=0.996669646468957\ldots\)\\

\addlinespace
\multicolumn{5}{@{}l}{\bfseries V. Criticité et mode \(30\)}\\*
phases \(3m-1\)
(\cref{eq:pi-cancellation,eq:chaos-closed,eq:899}) &
\(u_0=1-\cos(37y)\sin(36y)/(12\sin3y)\),
\(899=30^2-1\) &
Normalisation quadratique et annulation de \(\pi\) avant l'itération &
La constante de criticité se dérive d'un profil clos dont le mode \(30\)
sépare également \(90\) de \(120\) &
\(u_0(x)=x^2-(715769/2424603)x^4+O(x^6)\)\\

suite superstable
(\cref{eq:negative-schwarzian,eq:renormalization}) &
\(\mathcal R(f)(x)=-a^{-1}f(f(-ax))\) &
Unimodalité, schwarzienne négative et combinatoire de doublement &
Les quotients finis certifient l'entrée dans la classe universelle ;
l'hyperbolicité du point fixe reste une obligation séparée &
\(\delta_8=4.66921218915\ldots\),
\(\alpha_8=2.50296847988\ldots\)\\

\addlinespace
\multicolumn{5}{@{}l}{\bfseries VI. Spectre, arithmétique et cylindres}\\*
quart de torsion
(\cref{eq:J-square,eq:catalan}) &
\(J^2=-I\), \(G=L(2,\chi_{-4})\) &
Orientation complexe et multiplicativité du caractère &
Catalan est le deuxième moment de l'opérateur de quart de torsion, et non un
nombre décimal associé par ressemblance &
\(G=0.915965594177219\ldots\)\\

produit CRT
(\cref{eq:chi12,eq:L1chi12,eq:euler-kronecker}) &
\(\chi_{12}=\chi_{-3}\chi_{-4}\),
\(\gamma_{\mathbb Q(\sqrt3)}=\gamma+L'/L(1,\chi_{12})\) &
Orientation triadique, quart de torsion et terme fini de Dedekind &
Le réseau \(3/4/12\) détermine un corps réel et sa constante de bord &
\(L(1)=0.7603459963\ldots\),
\(\gamma_{\mathbb Q(\sqrt3)}=1.05396560825\ldots\)\\

zêta triadique renormalisée
(\cref{eq:bendersky,eq:odd-zeta-bendersky}) &
\(A_k=\exp(H_kB_{k+1}/(k+1)-\mathcal Z_3'(-k))\) &
Niveau de dérivée zêta et normalisation de Bernoulli &
Bendersky--Glaisher organise \(\zeta(2n+1)\) comme une famille graduée de
déterminations &
\(A_0=\sqrt{2\pi}\), \(A_1=A_{\rm GK}\), famille \(A_k\)\\

cylindres cofinaux
(\cref{eq:khinchin,eq:levy,eq:gauss-entropy,eq:lochs}) &
\(L=\pi^2/(12\log2)\), \(h_G=2L\),
\(\Lambda_b=\log b/h_G\) &
Mesure de Gauss, croissance des dénominateurs et rythme de codage &
Khinchin, Lévy et Lochs sont des observables distinctes du même système
transporté &
\(K=2.6854520011\ldots\), \(L=1.1865691104\ldots\),
\(\Lambda_{10}=0.9702701144\ldots\)\\

opérateur de transfert de Gauss
(\cref{eq:gkw-operator}) &
\(\mathcal L_Gf(x)=\sum_{n\ge1}(n+x)^{-2}
f((n+x)^{-1})\) &
Masse, mesure invariante et complémentaire de moyenne nulle &
La valeur propre sous-dominante mesure la relaxation des cylindres ; elle n'est pas
sélectionnée à partir d'une valeur décimale tabulée &
\(\lambda_{\rm GKW}\) : intervalle certifié à construire\\

cycles irréductibles et lacunes
(\cref{eq:meissel-mertens,eq:twin-prime-constant,eq:vacancy-gamma,eq:vacancy-stieltjes}) &
\(B_1=\lim(\sum_{p\le x}p^{-1}-\log\log x)\),
\(\gamma_{\rm vac,j}^+=\varepsilon_{\rm vac}\gamma_j+\cdots\) &
Factorisation première, discrépance bornée et holonomie de lacune &
Somme locale, densité multiplicative et termes finis conservent des
constructions antécédentes distinctes &
\(B_1=0.2614972128\ldots\), \(C_2=0.6601618158\ldots\),
\(\varepsilon_{\rm vac}=0.04575749056\ldots\)\\

relèvement \(r^2=2\pmod{3^k}\)
(\cref{eq:pell-order}) &
\(u=3+2r\), \(\operatorname{ord}(u)=4\cdot3^{k-1}\) &
Opérateur de quart de torsion compatible à toute profondeur nonadique &
L'unité de Pell conserve simultanément l'orientation finie et la mémoire
profinie &
\(3+2\sqrt2=5.82842712474619\ldots\)\\

\addlinespace
\multicolumn{5}{@{}l}{\bfseries VII. Changement de feuillet, mélange et électrofaible}\\*
\(C^*\mapsto-C^*\), \(6\to8\)
(\cref{eq:modular-sheet-parities,eq:modular-ckm-coordinates}) &
\((A,C^*)\mapsto(A,-C^*)\),
\(-(8-6)/3=-2/3\) &
Parité de feuillet et défaut orienté hexade--octade &
Barbero et CKM sont des réalisations distinctes du changement d'état, et non
des définitions rétrospectives mutuelles &
\(A,C^*\) retrouvés rationnellement à partir des trois angles\\

matrice de coordonnées \(M\)
(\cref{thm:modular-ckm-reflection,eq:modular-Rckm-matrix}) &
\(R_{\rm CKM}=M\operatorname{diag}(1,-1,1)M^{-1}\) &
Involution rationnelle, \(R^2=I\), \(\det R=-1\) &
L'espace de mélange conserve exactement la coordonnée qui a subi
l'inversion de feuillet ;
la phase CP reste fixe dans cette involution &
Matrice rationnelle exacte, \(\det M=-3857/6\)\\

opérateur angulaire \(T\)
(\cref{thm:ew-weinberg,eq:ew-sin2thetaW}) &
\(D_A=\det T=e^{-2A_{\rm rad}}\),
\(\sin^2\theta_W=1-D_A\) &
Déterminant pair et adjacence des routes \(W/Z\) &
Le vide et Weinberg partagent un opérateur angulaire, mais déterminent des fonctions
spectrales différentes &
\(D_A=0.775129119247\ldots\),
\(\sin^2\theta_W=0.224870880753\ldots\)\\

carte de jauge dans l'approximation à l'arbre
(\cref{eq:ew-g,eq:ew-gprime,eq:ew-GF,cor:ew-Higgs-lambda}) &
\(g=2\sqrt{\pi\alpha/(1-D_A)}\),
\(G_F^{(0)}=\pi\alpha/[\sqrt2(1-D_A)m_W^2]\) &
Le même \(D_A\) conserve la séparation du canal neutre/chargé &
Les couplages, la constante de Fermi et l'autocouplage de Higgs
forment une composition typée ; les corrections radiatives ne sont pas incorporées
à la fonctionnelle interne &
\(G_F^{(0)}=1.1157162817\cdot10^{-5}\,\mathrm{GeV}^{-2}\),
\(\lambda_H=0.123847911548\ldots\)\\

fonctionnelles \((\alpha,\varphi,1000,55)\)
(\cref{eq:atlas-anom-mu}) &
\(a_\mu=\alpha/(2\pi)+\alpha(\varphi-1)/1000+\alpha^2/(55000)\) &
Somme de trois contributions à coefficients figés &
Contrôle de composition leptonique ; ne remplace pas la série radiative complète
de QED &
\(a_\mu=0.001165920713\ldots\)\\

\addlinespace
\multicolumn{5}{@{}l}{\bfseries VIII. Cosmologie et contrôles de portée}\\*
rayon \(R_{108}=\ell_P\)
(\cref{eq:cosmo-ds-H-Lambda,eq:cosmo-ds-temperature,eq:cosmo-ds-entropy}) &
\(H_{108}=\pi/(54t_0)\),
\(\Lambda_{108}=\pi^2/(972\ell_0^2)\) &
Un seul rayon sous la réalisation de Sitter &
Expansion, courbure, température et entropie sont quatre grandeurs d'une
même géométrie conditionnelle &
\(T_{\rm dS}/\Theta_{\rm clk}=\log3/(2\pi)\),
\(S_{\rm dS}/k_B=\pi\)\\

échelle de Perron
(\cref{eq:cosmo-perron-a,eq:cosmo-minus-six,eq:cosmo-EC-scaling}) &
\(a_n/a_0=\varphi^{n/3}\),
\(M_n=(a_n/a_0)^{-6}\) &
Mode stable \(\varphi^{-2n}\) et dilution quadratique du spin &
L'exposant \(-6\) représente la mémoire active dans la loi cosmologique ;
l'amplitude d'Einstein--Cartan exige son transducteur &
\(M_na_n^6=a_0^6\),
\(H_\varphi t_0=0.1604039416865\ldots\)\\

contrôles de composition
(\cref{eq:atlas-bohr,eq:atlas-born}) &
\(a_0=\hbar/(m_ec\alpha)\),
\(E_*=h_{\rm int}c_{\rm int}/(e_{\rm int}\ell_0^2)\) &
Homogénéité dimensionnelle et non-alimentation par le nom &
Bohr se compose à partir de grandeurs définies hors du champ de cette
monographie ; Born--Infeld exige encore
une action déterminantale qui sélectionne son échelle &
Sections exactes \(a_0,E_*\), sans promouvoir de nouveau sélecteur\\

\end{longtable}
\endgroup

\begin{proof}
Chaque ligne s'obtient en substituant exclusivement des objets déjà construits dans
l'équation citée dans sa première ou sa deuxième colonne. Dans le vide, le calcul
est spectral : le théorème fonctionnel pour \(V\) produit \(r_\pm\), et les quatre
constitutives sont des fonctions de ces deux valeurs propres. Dans la réalisation
thermique, la structure chronologique \(108\) et l'entropie uniforme d'un
trit produisent le produit
\(k_B\Theta_{\rm clk}\), après quoi Planck, Wien et le rayonnement sont des
compositions. Dans la gravité, la condition circulaire fixe
\(\theta_{108}\), tandis que \(P_5\) et la réduction TT fixent le type du
support. Dans la réalisation modulaire, le système linéaire de déterminant \(30\)
reconstruit les deux canaux et la décomposition spectrale de l'état réduit
produit l'entropie et la première loi. Les constructions de criticité,
de caractères, de Gauss et de
Pell sont prouvées dans leurs opérateurs respectifs. La réflexion CKM s'obtient
par conjugaison d'une involution diagonale, et le secteur électrofaible par
la fonction déterminant de l'opérateur angulaire. Enfin, les deux
réalisations cosmologiques emploient, respectivement, un rayon périodicité temporelle de (108) phases et la
valeur propre stable
de Perron. À aucune étape, la valeur décimale de la cinquième colonne n'intervient pour
choisir l'antécédent. Cela prouve la terminalité et la confluence déclarées.
\end{proof}

\section{Relations structurelles principales de confluence généalogique}

Le corollaire précédent ne se limite pas à énumérer des lignes : il détermine quels
opérateurs partagent un antécédent et le point exact où leurs domaines se
séparent. Les cinq configurations suivantes contiennent les relations
causales de plus grande densité du volume.

\subsection{Mode \(30\) : complétion harmonique, criticité et mémoire}

Le même entier apparaît dans trois constructions exactes :
\begin{equation}
\begin{aligned}
90&=3\cdot30, &120&=4\cdot30,
&\frac{1-q^{90}}{1-q^{120}}
 &=\frac{1+s+s^2}{1+s+s^2+s^3},\qquad s=q^{30},\\
899&=30^2-1=29\cdot31,
&&&D&=90N_{90}+120N_{120}.
\end{aligned}
\label{eq:atlas-mode30-confluence}
\end{equation}
La première ligne est la complétion \(3\to4\) de l'opérateur du vide ; la
deuxième normalise le germe dodécaphasique sans introduire \(\pi\) dans
l'itération ; la troisième conserve le nombre d'occupations arrivées par chaque canal
modulaire. L'invariant commun est la partition du mode \(30\) ; les
codomaines sont différents. C'est pourquoi le croisement est structurel, et non une
identification de \(V\), \(u_0\) et \(\HHM\).

\paragraph{Identités de complétion du mode \(30\), inversibilité et condition d'incompatibilité.}
Les trois identités de \cref{eq:atlas-mode30-confluence} sont exactes. La
confluence est invalidée si \(90/120\) cessent d'être des multiples consécutifs du
même mode, si \(899\ne30^2-1\), ou si l'inversion
\((N,D)\mapsto(N_{90},N_{120})\) perd son déterminant \(30\). Les différences
de spectre ne la réfutent pas : ces différences sont précisément celles qui
maintiennent les trois réalisations typées.

\subsection{\(\log3\) : invariant résiduel, information et température}

La confluence informationnelle se résume dans la chaîne
\begin{equation}
S_{\TRIT}=k_B\log3,\qquad
E_P=k_B\Theta_{\rm clk}\log3,\qquad
I_{\rm or}=36\log3-S(\rho_A).
\label{eq:atlas-log3-confluence}
\end{equation}
Le premier membre mesure une décision ternaire uniforme ; le deuxième détermine
sa grandeur énergétique dans la structure chronologique \(108\) ; le troisième compare la
capacité de \(36\) trits à l'entropie de l'état réduit de bord.
Le scalaire est conservé,
mais l'unité, la température et l'état réduit restent typologiquement
séparés. Cette
séparation permet précisément de lire l'identité
\(E_P=k_B\Theta_{\rm clk}\log3\) : la réalisation de Planck apporte le quantum
d'énergie, la réalisation TRIT son contenu informationnel et la structure
chronologique le transducteur.

\paragraph{Conservation de \(\log 3\), séparation dimensionnelle et condition d'incompatibilité.}
La triple lecture est une composition exacte de
\cref{eq:log3-triple,eq:planck-trit,eq:modular-oriented-information}. Elle est
invalidée si l'on remplace \(\log3\) par une entropie binaire, si
\(\Theta_{\rm clk}\) est fixé à partir du kelvin avant l'identité, ou si
\(S(\rho_A)\) est calculé avec un état distinct de celui obtenu par trace partielle.

\subsection{\texorpdfstring{\(D_A\)}{Déterminant angulaire} : déterminant des feuillets et relation angulaire de Weinberg}

L'opérateur angulaire \(T=q_+P_++q_-P_-\) possède deux familles fonctionnelles :
\begin{equation}
F(T)^2\longrightarrow
(\widehat\varepsilon,\widehat\mu,\widehat Z,\widehat c),
\qquad
\det T=q_+q_-=D_A\longrightarrow
\sin^2\theta_W=1-D_A.
\label{eq:atlas-DA-confluence}
\end{equation}
La première conserve le spectre complet de la réponse du vide ; la
seconde écarte l'orientation individuelle des feuillets et ne conserve que leur
produit. D'où le fait que \(D_A\) soit pair sous \(C^*\mapsto-C^*\) et puisse
déterminer le quotient \(W/Z\), tandis que l'impédance conserve le quotient
des valeurs propres. La relation contient plus d'information qu'une
coïncidence angulaire : c'est
une bifurcation fonctionnelle exacte du même opérateur antécédent.

\paragraph{Exactitude du déterminant angulaire et obstructions de l'interface électrofaible.}
Le côté spectral et l'identité \(D_A=e^{-2A_{\rm rad}}\) sont exacts dans la
carte déclarée. La réalisation de jauge conserve son statut d'interface jusqu'à
la construction de l'entrelaceur de bases. Un quotient de routes non adjacentes, une
dépendance impaire de \(C^*\) ou \((m_W/m_Z)^2\ne D_A\) réfutent cette interface sans
affecter l'opérateur du vide.

\subsection{Secteurs \(90\) et \(120\) : vide, aire et dynamique modulaire}

Les canaux \(90\) et \(120\) ont trois fonctions non interchangeables :
ce sont des exposants de propagation dans \(V\), des multiplicités de l'incidence
\(6\to8\) et des poids spectraux dans \(\HHM\). Leur rapport \(4/3\) réapparaît dans
les températures modulaires,
\[
\frac{120}{90}=\frac43=\frac{\beta_{120}}{\beta_{90}},
\]
mais l'opérateur du vide n'est pas le hamiltonien modulaire et une température
modulaire n'est pas une perméabilité. La confluence consiste en ce que la même
généalogie conserve trois grandeurs : propagation, occupation aréale et
mémoire de canal.

\paragraph{Reconstruction finie des canaux \(90/120\) et hypothèses de la limite modulaire.}
La reconstruction finie et le rapport \(4/3\) sont exacts. La classification
de type III exige la persistance infinie des deux canaux. La disparition de l'un
d'eux à partir d'un certain niveau, la perte de compatibilité des états ou l'obtention
d'un groupe des rapports distinct de
\(\{\lambda_\partial^n:n\in\Z\}\) réfuteraient cette classification, et non
l'inversion finie.

\subsection{Paramètre de Barbero--Immirzi, entropie et CKM comme réalisations du passage \(6\to8\)}

L'interface de confluence complète s'écrit sans transformer une
réalisation en cause des autres :
\begin{equation}
(6\to8,\,90/120,\,C^*\mapsto-C^*)\longrightarrow
\begin{cases}
\gammaH\ell_P^2\longrightarrow A_{\rm phys},\\
\rho_A\longrightarrow(-\log\rho_A,S,I_{\rm or},\delta S),\\
C^*\mapsto-C^*\longrightarrow
R_{\rm CKM}=M\operatorname{diag}(1,-1,1)M^{-1}.
\end{cases}
\label{eq:atlas-bic-entropy-ckm}
\end{equation}
La première réalisation fixe le quantum géométrique ; la deuxième mesure
l'information de l'état réduit et sa réponse aréale ; la troisième enregistre
l'orientation du feuillet dans
l'espace de mélange. En parallèle,
\(E_P=k_B\Theta_{\rm clk}\log3\) fixe le quantum informationnel d'énergie.
Ainsi, gravité, entropie et information sont mises en équations au moyen de
\((\ell_P^2,E_P)\), sans identifier \(\gammaH\), \(k_B\) ou les angles CKM.

\paragraph{Indépendance causale de Barbero--Immirzi, de l'information modulaire et du mélange CKM.}
La sortie \(\gammaH\) est un résultat antécédent démontré de sa dérivation correspondante ;
l'opérateur d'aire, la purification, la première loi locale et la réflexion CKM
sont exacts dans cette monographie. La classification modulaire infinie et la lecture
physique d'horizon conservent leurs hypothèses. Le diagramme est invalidé si
l'on utilise CKM pour recalibrer \(\gammaH\), si l'on assimile une hexade/octade
d'incidence à des faces/sommets sans entrelaceur, ou si
\(R_{\rm CKM}^2\ne I\) ou \(\det R_{\rm CKM}\ne-1\).

\section{Indépendance théorématique du corollaire de confluence métrologique}

La \cref{tab:metrological-confluence} ne renvoie pas à un tableau extérieur pour
compléter l'une quelconque de ses cinq colonnes. Chaque antécédent est défini dans le
corps ; chaque équation est démontrée ou typée dans le localisateur cité ;
chaque invariant déclare ce qui est conservé et aussi ce qui est écarté ; chaque
signification empêche l'identification par la valeur décimale ; et chaque valeur terminale
est obtenue après cette chaîne. Les statuts et critères de réfutation restent dans les
fiches détaillées qui suivent, où ils sont développés sans la compression inévitable
d'un tableau synoptique.

En conséquence, le tableau ne remplace pas le développement : il le complète. Il peut se lire
de manière synoptique pour reconnaître la confluence, mais sa validité procède des
théorèmes et des preuves du corps, et chaque ligne dispose ensuite d'une
fiche autonome avec antécédent, équation, signification, valeur, statut et
contrôle négatif.

\section{Action, vide et constantes quantiques--électriques}

\subsection*{Sections bidirectionnelles de l'action }
\textbf{Antécédent.} Deux réalisations orientées d'une même droite
d'action. \textbf{Équation.}
\[
R_{\rm act}=\frac{\hbar_{\rm pre}}{\hbar_{\rm ret}}
=\frac{H_5}{\eta_{\rm ret}},
\qquad \xi_{\rm act}=\tfrac12\log R_{\rm act}.
\]
\textbf{Structure et signification.} La décennie commune fixe le type ; le quotient
conserve l'orientation. \textbf{Valeur.} \(R_{\rm act}\) et
\(\xi_{\rm act}\) dans la carte interne, sans comparer des mantisses dimensionnelles à des
résidus. \textbf{Statut.} \emph{Théorème interne exact}.
\textbf{Critère de réfutation.} Utiliser la décennie commune comme correcteur relatif ou mélanger
les deux sections.

\subsection*{Opérateur constitutif du vide : spectre positif et reconstruction conjointe de \(\varepsilon,\mu,Z,c\)}
\textbf{Antécédent.} L'opérateur bicouche \(T=q_+P_++q_-P_-\) et la généalogie
\(30\to90/120\). \textbf{Équation.}
\[
V=(I-T^{90})(I-T^{120})^{-1},
\qquad
\frac{1-q^{90}}{1-q^{120}}
=\frac{1+s+s^2}{1+s+s^2+s^3},\quad s=q^{30}.
\]
Si \(\Spec V=\{r_+,r_-\}\),
\[
\widehat\varepsilon=r_+^2,
\quad\widehat\mu=r_-^2,
\quad\widehat Z=r_-/r_+,
\quad\widehat c=(r_+r_-)^{-1}.
\]
\textbf{Structure et signification.} Les quatre constitutives sont des fonctions
d'un seul opérateur positif ; chaque feuillet complète trois segments par un quatrième segment.
L'amplification \(V_m=V\otimes I_{9^m}\) conserve les fonctionnelles sous
l'espérance traciale.
\textbf{Valeur.}
\begin{align*}
\widehat\varepsilon&=0.9999992570966367027604416155\ldots,\quad
\widehat\mu=0.9994483504793269350899815559\ldots,\\
\widehat Z&=0.9997245085389372154555543466\ldots,\quad
\widehat c=1.0002763104861575261063862689\ldots.
\end{align*}
\textbf{Statut.} \emph{Théorème interne exact} ; amplification nonadique :
\emph{Formalisation nouvelle exacte}. \textbf{Critère de réfutation.} Perte de
positivité, résidu dans l'identité \(3\to4\), ou variation des
invariants sous l'espérance traciale.

\subsection*{Sections constitutives de \(\varepsilon\), \(\mu\), \(Z\) et \(c\) dans les cartes orientées}
\textbf{Antécédent.} Les quatre coefficients précédents et le cadre
\((u_S,u_C,u_T,u_Q)\). \textbf{Équation.}
\[
Z=\widehat Zu_Su_Q^{-2},\quad c=\widehat cu_C,\quad
\mu=\widehat\mu u_Su_C^{-1}u_Q^{-2},\quad
\varepsilon=\widehat\varepsilon u_S^{-1}u_C^{-1}u_Q^2.
\]
\textbf{Structure et signification.} La permittivité et la perméabilité ne
reçoivent pas d'unités après un calcul adimensionnel : ce sont des sections de droites
constitutives. \textbf{Valeur.} Les quatre sections précédentes, avec
\(\mu\varepsilon=c^{-2}\) et \(\mu/\varepsilon=Z^2\).
\textbf{Statut.} \emph{Théorème interne exact}. \textbf{Critère de réfutation.} Introduire
\(c,\mu_0,\varepsilon_0\) SI avant l'opérateur ou violer l'homogénéité.

\subsection*{Dérivation conjointe de la charge, de la résistance de von Klitzing, de la conductance, du flux et de la constante de Josephson}
\textbf{Antécédent.} \((Z_{\rm vac},h_a,\alpha_{\HMT})\).
\textbf{Équation.}
\[
e_a=\sqrt{\frac{2\alpha_{\HMT}h_a}{Z_{\rm vac}}},\quad
R_K=\frac{Z_{\rm vac}}{2\alpha_{\HMT}},\quad
G_0=\frac{4\alpha_{\HMT}}{Z_{\rm vac}},
\]
\[
\Phi_{0,a}=\sqrt{\frac{h_aZ_{\rm vac}}{8\alpha_{\HMT}}},
\qquad K_{J,a}=\Phi_{0,a}^{-1}.
\]
\textbf{Structure et signification.} \(R_K\) et \(G_0\) sont des invariants
bidirectionnels ; \(\Phi_0\) et \(K_J\) conservent l'orientation de l'action avec des
exposants opposés. \textbf{Valeur.} Les expressions exactes dans leurs droites ;
\(R_KG_0=2\) et \(K_J\Phi_0=1\). \textbf{Statut.}
\emph{Théorème interne exact}. \textbf{Critère de réfutation.} Utiliser la charge SI comme sélecteur ou
échouer à satisfaire \(\alpha=Ze^2/(2h)\).

\section{Constantes thermiques et rayonnement}

\subsection*{Relation de Boltzmann et contenu entropique TRIT }
\textbf{Antécédent.} Décision ternaire uniforme, section d'action et
structure chronologique \(108\). \textbf{Équation.}
\[
k_B\Theta_{\rm clk}\log3=\frac{h_{\rm int}}{108t_0}
=\frac{2\pi\hbar_{\rm int}}{108t_0},
\qquad
\log3=\frac{S_{\TRIT}}{k_B}=\frac{E_P}{k_B\Theta_{\rm clk}}.
\]
\textbf{Structure et signification.} \(\log3\) est simultanément un défaut
spectral, le contenu d'une décision et un rapport thermique ; les opérateurs
restent
distincts. \textbf{Valeur.} \(\log3=1.0986122886681098\ldots\) et la section
énergie--température précédente. \textbf{Statut.} \emph{Confluence
exacte}. \textbf{Critère de réfutation.} Séparer arbitrairement \(k_B\) de
\(\Theta_{\rm clk}\) ou inférer l'identité des domaines à partir de l'égalité du scalaire.

\subsection*{Loi de déplacement de Wien }
\textbf{Antécédent.} Le spectre de Planck évalué en
\(\Theta_{\rm clk}\). \textbf{Équation.}
\[
x_5=5+W_0(-5e^{-5}),\quad
\frac{\lambda_{\max}}{\ell_0}=\frac{108\log3}{x_5},
\]
\[
x_3=3+W_0(-3e^{-3}),\quad
\nu_{\max}t_0=\frac{x_3}{108\log3}.
\]
\textbf{Structure et signification.} Longueur d'onde et fréquence sont deux
paramétrisations spectrales avec des jacobiens distincts. \textbf{Valeur.}
\(\lambda_{\max}/\ell_0=23.896756779\ldots\) et
\(\nu_{\max}t_0=0.0237794888\ldots\). \textbf{Statut.}
\emph{Théorème interne exact}. \textbf{Critère de réfutation.} Confondre \(x_5\) avec \(x_3\) ou
introduire le kelvin avant le produit thermique.

\subsection*{Loi de Stefan--Boltzmann et gaz de photons }
\textbf{Antécédent.} La même échelle thermique, le continu bosonique et la
fonctionnelle \(\zeta(3)\). \textbf{Équation.}
\[
\mathcal F(\Theta_{\rm clk})
=\frac{2\pi^5}{15\,108^4(\log3)^4}
\frac{h_{\rm int}}{\ell_0^2t_0^2},
\]
\[
n_\gamma\ell_P^3=\frac{2\zeta(3)}{\pi^2(\log3)^3},
\qquad
\frac{u_\gamma\ell_P^3}{E_P}=\frac{\pi^2}{15(\log3)^4}.
\]
\textbf{Structure et signification.} \(\zeta(3)\) compte l'occupation bosonique ;
\(\log3\) fixe le contenu informationnel TRIT par cellule. \textbf{Valeur.}
Les coefficients
exacts précédents. \textbf{Statut.} \emph{Composition exacte avec
structure de départ explicite}. \textbf{Critère de réfutation.} Perte de facteurs
\(2\pi\), dimensions de flux incorrectes ou fusion des provenances mathématiques de
\(\zeta(3)\) et \(\log3\).

\section{Criticité et constantes dynamiques}

\subsection*{Profil dodécaphasique et normalisation par \(899\)}
\noindent\textit{Résultats directeurs : , et
.}\par
\textbf{Antécédent.} Douze phases \(3m-1\) et mode élémentaire \(30\).
\textbf{Équation.} Avec \(y=x/\sqrt{899}\),
\[
u_0(x)=1-\frac1{12}\sum_{m=1}^{12}
\cos\!\bigl(2(3m-1)y\bigr)
=1-\frac{\cos37y\sin36y}{12\sin3y},
\]
\[
899=30^2-1=29\cdot31,\qquad
u_0(x)=x^2-\frac{715769}{2424603}x^4+O(x^6).
\]
\textbf{Structure et signification.} La phase \(\pi\) s'annule avant
l'itération ; le profil résultant est analytique, unimodal quadratique et de
schwarzienne négative sur \([-1,1]\). \textbf{Valeur.} Profil et coefficients
rationnels exacts. \textbf{Statut.} \emph{Théorème interne exact}.
\textbf{Critère de réfutation.} Un deuxième point critique, une schwarzienne non négative ou un
résidu dans l'identité trigonométrique.

\subsection*{Approximants superstables de Feigenbaum et certificat hyperbolique d'universalité}
\textbf{Antécédent.} La suite superstable du germe précédent.
\textbf{Équation.}
\[
\mathcal R(f)(x)=-a(f)^{-1}f\bigl(f(-a(f)x)\bigr),
\qquad a(f)=-f(1).
\]
\textbf{Structure et signification.} Les quotients finis prouvent que la
famille entre dans la route appropriée ; l'universalité exige le point fixe et
son hyperbolicité. \textbf{Valeur.}
\(\delta_8=4.66921218915\ldots\),
\(\alpha_8=2.50296847988\ldots\). \textbf{Statut.} Approximants :
\emph{Certificat fini} ; limite universelle : \emph{Programme avec
critère de réfutation}. \textbf{Critère de réfutation.} Plus d'une valeur propre instable, queue non bornée
ou transversalité nulle.

\section{Constantes spectrales, arithmétiques et de lacune}

\subsection*{Structure complexe interne et constante de Catalan }
\textbf{Antécédent.} \(J^2=-I\) et le caractère
\(\chi_J(n)=i^{n-1}=\chi_{-4}(n)\) sur les impairs. \textbf{Équation.}
\[
G=L(2,\chi_J)=\sum_{n\ge0}\frac{(-1)^n}{(2n+1)^2}.
\]
\textbf{Structure et signification.} Catalan est le moment quadratique du
quart de torsion. \textbf{Valeur.} \(G=0.915965594177219\ldots\).
\textbf{Statut.} \emph{Théorème interne exact}. \textbf{Critère de réfutation.} Absence de
multiplicativité ou désaccord avec \(\chi_{-4}\).

\subsection*{Caractère dodécaphasique et corps quadratique réel }
\textbf{Antécédent.} Produit CRT
\(\chi_{12}=\chi_{-3}\chi_{-4}\). \textbf{Équation.}
\[
L(1,\chi_{12})=\frac{\log(2+\sqrt3)}{\sqrt3},
\qquad L(2,\chi_{12})=\frac{\pi^2}{6\sqrt3},
\]
\[
\gamma_{\mathbb Q(\sqrt3)}
=\gamma+\frac{L'(1,\chi_{12})}{L(1,\chi_{12})}.
\]
\textbf{Structure et signification.} L'orientation triadique et le quart de
torsion produisent le caractère réel de \(\mathbb Q(\sqrt3)\) ;
Euler--Kronecker est
le terme fini de sa zêta de Dedekind. \textbf{Valeur.}
\(0.7603459963009463\ldots\), \(0.9497031262940094\ldots\) et
\(1.05396560824848258\ldots\). \textbf{Statut.} \emph{Théorème interne exact} ;
l'identité composée conserve les hypothèses et les normalisations déclarées.
\textbf{Critère de réfutation.} Motif de résidus distinct de \((+,-,-,+)\) ou échec de
la factorisation de Dedekind.

\subsection*{Famille Bendersky--Glaisher }
\textbf{Antécédent.} Fonctionnelle zêta triadique renormalisée.
\textbf{Équation.}
\[
A_k=\exp\!\left(\frac{H_kB_{k+1}}{k+1}-\mathcal Z_3'(-k)\right),
\quad
\zeta(2n+1)=(-1)^{n+1}\frac{2(2\pi)^{2n}}{(2n)!}\log A_{2n}.
\]
\textbf{Structure et signification.} Apéry et les valeurs impaires de zêta sont des
niveaux d'une tour, et non des coïncidences isolées. \textbf{Valeur.}
\(A_0=\sqrt{2\pi}\), \(A_1=A_{\rm GK}\) et la famille exacte précédente.
\textbf{Statut.} \emph{Théorème interne exact}. \textbf{Critère de réfutation.} Échec dans
\(A_0,A_1\) ou dans l'identité de \(\zeta(2n+1)\).

\subsection*{Système dynamique de Gauss : \(K\), \(L\), \(\lambda_{\mathrm{GKW}}\) et cylindres transportés}
\textbf{Antécédent.} Évaluation des cylindres HMT et mesure
\(d\mu_G=dx/[\log2(1+x)]\). \textbf{Équation.}
\[
\log K=\int\log a_1\,d\mu_G,\quad
L=-\int\log x\,d\mu_G=\frac{\pi^2}{12\log2},
\]
\[
h_G=2L=\frac{\pi^2}{6\log2},\qquad
\Lambda_b=\frac{6\log2\log b}{\pi^2}.
\]
\textbf{Structure et signification.} Khinchin mesure les chiffres ; Lévy mesure la
croissance des dénominateurs ; Lochs compare les rythmes de codage.
\textbf{Valeur.}
\begin{align*}
K&=2.685452001065306\ldots,
&L&=1.18656911041562545\ldots,\\
h_G&=2.37313822083125091\ldots,
&\Lambda_{10}&=0.97027011439203393\ldots,\\
&&\Lambda_{729}&=2.77761896637430578\ldots.
\end{align*}
\textbf{Statut.} Cofinalité et identités : \emph{Exact dans le système
transporté} ; constante GKW : \emph{Programme ouvert avec critère de réfutation}.
\textbf{Critère de réfutation.} Non-cofinalité des cylindres, mesure non invariante ou absence
de trou spectral dans l'espace déclaré.

\subsection*{Cycles arithmétiques de Meissel--Mertens et séries régularisées de nombres premiers}
\textbf{Antécédent.} Sélecteur de cycles irréductibles et factorisation
exacte. \textbf{Équation.}
\[
B_1=\lim_{x\to\infty}\left(\sum_{p\le x}\frac1p-\log\log x\right),
\qquad
C_2=\prod_{p>2}\frac{p(p-2)}{(p-1)^2}.
\]
\textbf{Structure et signification.} \(B_1\) est une somme locale régularisée ;
\(C_2\) une densité multiplicative de paires. \textbf{Valeur.}
\(B_1=0.2614972128476428\ldots\),
\(C_2=0.6601618158468696\ldots\). \textbf{Statut.}
\emph{Démontré avec structure de départ explicite} : sélecteur premier et
borne de queue déclarés. \textbf{Critère de réfutation.} Utiliser des zéros ou des décimales pour choisir
les nombres premiers, ou affirmer la limite sans appliquer la borne de queue.

\subsection*{Lacunes, holonomie et constantes de Stieltjes }
\textbf{Antécédent.} Mot mécanique de lacunes et discrépance bornée.
\textbf{Équation.}
\[
\varepsilon_{\rm vac}=\log_{10}\frac{10}{9},\quad
\gamma_{\rm vac}^+=\lim_N\left(\sum_{n\le N}\frac{z_n}{n}
-\varepsilon_{\rm vac}\log N\right),
\]
\[
\rho_\varepsilon(k)=e^{2\pi ik\varepsilon_{\rm vac}},\qquad
\gamma_{\rm vac,j}^+=\varepsilon_{\rm vac}\gamma_j+
\sum_{n\ge1}\frac{(z_n-\varepsilon_{\rm vac})(\log n)^j}{n}.
\]
\textbf{Structure et signification.} Densité, terme fini et phase sont trois
types ; la tour prolonge le terme fini. \textbf{Valeur.}
\(\varepsilon_{\rm vac}=0.0457574905606751\ldots\) ; les autres sont déterminés
par leurs limites certifiées. \textbf{Statut.} \emph{Théorème interne exact avec
discrépance}. \textbf{Critère de réfutation.} Perte de la borne de discrépance ou
désaccord du niveau \(j=0\).

\subsection*{Relèvement nonadique de Pell }
\textbf{Antécédent.} \(r^2=2\pmod{3^k}\). \textbf{Équation.}
\[
u=3+2r,\qquad \operatorname{ord}(u)=4\cdot3^{k-1},
\qquad \varprojlim\langle u\rangle=C_4\times\mathbb Z_3.
\]
\textbf{Structure et signification.} Le quart de tour persiste à toute
profondeur ; la réalisation archimédienne est \((1+\sqrt2)^2\).
\textbf{Valeur.} \(3+2\sqrt2=5.82842712474619\ldots\).
\textbf{Statut.} \emph{Théorème interne exact}. \textbf{Critère de réfutation.} Ordre distinct
ou relèvements incompatibles.

\section{Gravité, aire et information modulaire}

\subsection*{Constante gravitationnelle et représentation tensorielle à cinq composantes }
\textbf{Antécédent.} Structure chronologique \(108\), action et projecteur
central de rang
cinq. \textbf{Équation.}
\[
G_a(\theta)=\theta\frac{c_{\rm int}^5t_0^2}{\hbar_a},
\qquad \theta_{108}=\left(\frac{54}{\pi}\right)^2,
\qquad \mathcal G_5=G_a(\theta_{108})P_5.
\]
\textbf{Structure et signification.} \(G\) est une section sélectionnée par la
clôture Compton--gravité ; \(P_5\) détermine un support à cinq composantes,
que la réduction TT ramène à deux hélicités sans masse.
\textbf{Valeur.} \(G_{\rm int}=\theta_{108}c^5t_0^2/\hbar\) et
\(\Spec\mathcal G_5=\{G\text{ avec multiplicité }5,
0\text{ avec multiplicité }7\}\).
\textbf{Statut.} Sélecteur : \emph{Théorème interne exact} ; réalisation du support :
\emph{Réalisation physique typée}. \textbf{Critère de réfutation.} Une autre racine positive,
\(\Tr P_5\ne5\), ou confondre cinq composantes avec cinq polarisations
physiques de GR sans masse.

\subsection*{Système de Planck, réciprocité \(G\leftrightarrow\hbar\) et invariance aréale}
\textbf{Antécédent.} Le produit bidirectionnel \(G_a\hbar_a\) et la périodicité temporelle de (108) phases.
\textbf{Équation.}
\[
\ell_P^2=\frac{G\hbar}{c^3}=\theta_{108}\ell_0^2,
\quad t_P=\frac{54}{\pi}t_0,
\quad E_P=\frac{\pi}{54}\frac{\hbar}{t_0},
\]
\[
F_P=\frac{c^4}{G},\qquad P_P=\frac{c^5}{G},\qquad
\rho_P=\frac{\hbar}{\theta_{108}^2c^5t_0^4}.
\]
\textbf{Structure et signification.} L'aire de Planck est invariante sous le
dédoublement de l'action et relie la gravité à l'échelle thermique.
\textbf{Valeur.} Les sections exactes précédentes.
\textbf{Statut.} \emph{Théorème interne exact}. \textbf{Critère de réfutation.} Combiner
\(G\) et \(\hbar\) de feuillets incompatibles ou échouer à satisfaire l'homogénéité.

\subsection*{Opérateur d'aire et hamiltonien modulaire : \(\widehat{\mathcal A}_{\rm HMT}\), \(H_{\rm mod}\) et canaux \(90/120\)}
\textbf{Antécédent.} Incidence \(6\to8\), canaux \(90/120\), sortie
\(\gammaH\) et \(\ell_P^2\). \textbf{Équation.}
\[
\frac{A_{\rm phys}}{\ell_P^2}=\gammaH a_0(N_{90}+N_{120}),
\qquad
\HHM=\Dmod(90N_{90}+120N_{120})+cI.
\]
Si \(N=N_{90}+N_{120}\) et
\(D=90N_{90}+120N_{120}\),
\[
N_{90}=4N-D/30,\qquad N_{120}=D/30-3N.
\]
\textbf{Structure et signification.} L'aire détermine l'occupation totale ; le
hamiltonien conserve la provenance du canal. Ensemble, ils reconstruisent le bord.
\textbf{Valeur.} \(\gammaH a_0=0.0016755940749224991\ldots\) dans la carte
déclarée. \textbf{Statut.} \emph{Théorème interne exact en coupe finie} ; on ne
redérive pas \(\gammaH\). \textbf{Critère de réfutation.} Commutateur non nul, déterminant
\(30\) perdu ou emploi de l'adjectif holographique sans transport au bord.

\subsection*{Entropie relative, purification thermochamp double et première loi modulaire aréale}
\textbf{Antécédent.} État réduit fidèle du bord et sa purification de
Schmidt. \textbf{Équation.}
\[
I_{\rm or}=36\log3-S(\rho_A),
\quad
|\mathrm{TFD}\rangle=\sum_n\sqrt{p_n}|n\rangle|n\rangle,
\]
\[
\delta S=\beta_{90}\delta\langle \mathcal A_{90}\rangle
+\beta_{120}\delta\langle \mathcal A_{120}\rangle,
\qquad \frac{\beta_{120}}{\beta_{90}}=\frac43.
\]
\textbf{Structure et signification.} La purification transforme l'entropie
réduite en intrication ; la première loi met en équations mémoire modulaire et aire.
\textbf{Valeur.} \(I_{\rm or}=0.0016691832338689407\ldots\),
\(\beta_{90}=5.97264854010709\ldots\),
\(\beta_{120}=7.96353138680946\ldots\).
\textbf{Statut.} \emph{Théorème interne exact local}. \textbf{Critère de réfutation.} Trace
partielle distincte de \(\rho_A\), \(I_{\rm or}<0\), ou utilisation pour des variations
finies sans le terme d'entropie relative.

\subsection*{Limite modulaire et paramètre de type III }
\textbf{Antécédent.} États locaux compatibles avec persistance infinie des
deux canaux. \textbf{Équation.}
\[
\HHM^{(m)}=-\log\rho_m,
\qquad
\lambda_\partial=\exp(-30\Dmod).
\]
\textbf{Structure et signification.} À l'infini vit un réseau local ou le
générateur modulaire GNS, et non une matrice de densité globale dans l'UHF.
\textbf{Valeur.}
\(\lambda_\partial=0.9966696464689573786\ldots\).
\textbf{Statut.} Réseau : \emph{Règle de type opératoriel} ; classification :
\emph{Théorème avec hypothèses de persistance}. \textbf{Critère de réfutation.} Absence
asymptotique d'un canal, état non produit ou rapport asymptotique distinct.

\section{Involution des feuillets, mélange et secteur électrofaible}

\subsection*{Involution des feuillets et transformation rationnelle CKM : unitarité, déterminant et caractère de Jarlskog}
\textbf{Antécédent.} Échange \(C^*\mapsto-C^*\) et donnée
\((h_6,o_8,t)=(6,8,3)\). \textbf{Équation.}
\[
R_{\rm CKM}=M\operatorname{diag}(1,-1,1)M^{-1},
\quad R_{\rm CKM}^2=I,\quad\det R_{\rm CKM}=-1,
\]
\[
-\frac{o_8-h_6}{t}=-\frac23,
\quad
A=\frac{1528\theta_{12}+3380\theta_{23}+801\theta_{13}}{3857},
\]
\[
C^*=6\frac{75\theta_{12}+176\theta_{23}-150\theta_{13}}{3857}.
\]
\textbf{Structure et signification.} Barbero et CKM sont des réalisations
fonctionnelles distinctes du changement \(6\to8\), adopté comme antécédent
commun ; l'involution induit une réflexion rationnelle de rang impair un.
\textbf{Valeur.} Matrice rationnelle exacte, \(\det M=-3857/6\) ; la phase CP
reste fixe sous cette involution, qui n'est pas une conjugaison CP.
\textbf{Statut.} \emph{Exact généalogique/interne}.
\textbf{Critère de réfutation.} \(R^2\ne I\), échec de \(RM=MD\), perte de
\(-2/3\), ou présentation de CKM comme cause rétrospective de \(\gammaH\).

\subsection*{Déterminant angulaire, paramètre de Weinberg et couplages électrofaibles}
\textbf{Antécédent.} Deux fonctions du même opérateur angulaire :
\(F(T)^2\) et \(\det T\) ; routes \(W/Z\) adjacentes.
\textbf{Équation.}
\[
D_A=\det T=e^{-2A_{\rm rad}},
\qquad \sin^2\theta_W^{\rm OS}=1-D_A,
\]
\[
g=2\sqrt{\frac{\pi\alpha}{1-D_A}},
\qquad g'=2\sqrt{\frac{\pi\alpha}{D_A}}.
\]
\textbf{Structure et signification.} Le vide et Weinberg sont des réalisations
spectrales distinctes d'un opérateur commun ; le déterminant est pair sous le
changement de feuillet.

\textbf{Valeur.} \(D_A=0.7751291192471564302\ldots\) et
\(\sin^2\theta_W=0.2248708807528435698\ldots\).

\textbf{Statut.} Relation angulaire : \emph{Théorème interne exact dans la carte} ;
réalisation de jauge : \emph{Réalisation conditionnelle}.
\textbf{Critère de réfutation.} Routes non adjacentes, quotient de masses
distinct de \(D_A\), ou schéma de jauge non
déclaré.

\subsection*{Constante de Fermi et couplage de Higgs }
\textbf{Antécédent.} \(\alpha_{\HMT}\), \(D_A\), routes \(W/H\) et carte
électrofaible à l'arbre. \textbf{Équation.}
\[
G_F^{(0)}=\frac{\pi\alpha}{\sqrt2(1-D_A)m_W^2},
\]
\[
\lambda_H=\frac{\pi\alpha}{2(1-D_A)}
\exp\!\left(6A_{\rm rad}+\frac{10}{3}C^*_{\rm rad}\right),
\qquad G_F^{(0)}m_H^2=\sqrt2\lambda_H.
\]
\textbf{Structure et signification.} Fermi est la limite effective de basse
énergie du canal chargé ; Higgs conserve l'orientation que le
déterminant de jauge élimine.
\textbf{Valeur.} \(G_F^{(0)}=1.1157162817\times10^{-5}\,\mathrm{GeV}^{-2}\)
dans la carte déclarée et \(\lambda_H=0.123847911548\ldots\).
\textbf{Statut.} \emph{Interface électrofaible à l'arbre}.
\textbf{Critère de réfutation.} Utiliser \(G_F\) observé pour choisir \(\Delta r\), omettre
le schéma/l'échelle ou sélectionner \(\lambda_H\) à partir de la masse observée.

\subsection*{Anomalie magnétique leptonique }
\textbf{Antécédent.} \((\alpha,\varphi,1000,54+1)\) avec des coefficients
figés avant la comparaison. \textbf{Équation.}
\begin{equation}
a_\mu=\frac{\alpha}{2\pi}
+\frac{\alpha(\varphi-1)}{1000}
+\frac{\alpha^2}{1000(54+1)}.
\label{eq:atlas-anom-mu}
\end{equation}
\textbf{Structure et signification.} C'est une évaluation de trois fonctionnelles HMT,
et non une substitution à la série radiative complète de QED.
\textbf{Valeur.} \(a_\mu=0.001165920713\ldots\).
\textbf{Statut.} \emph{Évaluation exacte avec provenance locale ouverte}.
\textbf{Critère de réfutation.} Choisir \(1000\) ou \(55\) à partir de la valeur expérimentale.

\section{Détermination des paramètres cosmologiques par mémoire torsionnelle et dynamique de Perron}

\subsection*{Réalisation de de Sitter par la périodicité temporelle de \(108\) phases}
\textbf{Antécédent.} Sélecteur circulaire \(R_{108}=\ell_P\) et réalisation de
Sitter. \textbf{Équation.}
\[
H_{108}=\frac{\pi}{54t_0},\quad
\Lambda_{108}=\frac{\pi^2}{972\ell_0^2},\quad
\frac{T_{\rm dS}}{\Theta_{\rm clk}}=\frac{\log3}{2\pi},\quad
\frac{S_{\rm dS}}{k_B}=\pi.
\]
\textbf{Structure et signification.} Un même rayon détermine l'expansion,
la courbure, la température et l'entropie. \textbf{Valeur.} Le quadruplet exact
précédent. \textbf{Statut.} \emph{Réalisation physique conditionnelle}.
\textbf{Critère de réfutation.} Fond non de Sitter ou réalisateur PCH sélectionnant un autre
rayon.

\subsection*{Réalisation de Perron--Einstein--Cartan }
\textbf{Antécédent.} Échelle \(V_n=\varphi^n\), \(D=3\), mode stable
\(\varphi^{-2n}\). \textbf{Équation.}
\[
\frac{a_n}{a_0}=\varphi^{n/3},\qquad
M_n=\left(\frac{a_n}{a_0}\right)^{-6},\qquad
H_\varphi=\frac{\log\varphi}{3t_0}.
\]
\textbf{Structure et signification.} L'exposant \(-6\) est une mémoire active et
coïncide avec la dilution quadratique du spin d'Einstein--Cartan.
\textbf{Valeur.} \(H_\varphi t_0=0.1604039416865344825\ldots\) et la loi exacte
\(M_na_n^6=a_0^6\). \textbf{Statut.} Échelle :
\emph{Théorème interne exact} ; lecture du spin : \emph{Réalisation conditionnelle}.
\textbf{Critère de réfutation.} Exposant distinct, paramètre temporel non
additif ou amplitude physique
choisie à partir d'un rebond.

\section{Séparation typée entre les rayons de Bohr et de Born--Infeld et les échelles engendrées en interne}

\subsection*{Rayon de Bohr comme échelle électrodynamique dérivée de \(\alpha\)}
\textbf{Antécédent.} Action, vide, \(\alpha\) et masse électronique de son
antécédent. \textbf{Équation.}
\begin{equation}
a_0=\frac{\hbar}{m_ec\alpha}.
\label{eq:atlas-bohr}
\end{equation}
\textbf{Signification et valeur.} C'est une composition
descendante, et non un sélecteur indépendant de cette monographie.
\textbf{Statut.} \emph{Hors du champ}. \textbf{Critère de réfutation.} Recalculer
\(m_e\) ici ou utiliser \(a_0\) pour sélectionner \(\alpha\).

\subsection*{Échelle de Born--Infeld et récupération de la limite de Maxwell}
\textbf{Antécédent.} \((h_{\rm int},c_{\rm int},e_{\rm int},\ell_0)\).
\textbf{Équation.}
\begin{equation}
E_*=\frac{h_{\rm int}c_{\rm int}}{e_{\rm int}\ell_0^2}.
\label{eq:atlas-born}
\end{equation}
\textbf{Structure et signification.} L'expression construit une échelle de
champ ; elle ne construit pas encore une action de Born--Infeld.
\textbf{Valeur.} Section exacte \(E_*\), sans détermination d'un
paramètre physique.
\textbf{Statut.} \emph{Programme ouvert}. \textbf{Critère de réfutation.} Identifier
par le nom ou la valeur décimale sans action déterminantale ni sélecteur.

\begin{proposition}[Résultat principal]
Le résultat métrologique n'est pas la valeur décimale isolée. C'est la conservation
d'une généalogie à travers des opérateurs, des droites de grandeur et des cartes, jusqu'à
ce que la valeur réelle apparaisse à la dernière position de la fiche.
\end{proposition}
